
\documentclass{article}
\usepackage{colortbl}
\usepackage{makecell}
\usepackage{multirow}
\usepackage{supertabular}

\begin{document}

\newcounter{utterance}

\centering \large Interaction Transcript for game `bbh', experiment `salient\_translation\_error\_detection', episode 23 with reflexion{-}Qwen3.5{-}9B.
\vspace{24pt}

{ \footnotesize  \setcounter{utterance}{1}
\setlength{\tabcolsep}{0pt}
\begin{supertabular}{c@{$\;$}|p{.15\linewidth}@{}p{.15\linewidth}p{.15\linewidth}p{.15\linewidth}p{.15\linewidth}p{.15\linewidth}}
    \# & \multicolumn{2}{c}{Player} && \multicolumn{2}{c}{Game Master} \\
    \hline

    \theutterance \stepcounter{utterance}  
    & & & \multicolumn{4}{p{0.6\linewidth}}{
        \cellcolor[rgb]{0.9,0.9,0.9}{
            \makecell[{{p{\linewidth}}}]{
                \texttt{\tiny{[P1$\langle$GM]}}
                \texttt{You have attempted this same task before and did not succeed. These are the plans you wrote after those attempts.} \\
\texttt{Trial \#1: New plan: I incorrectly classified the shortening of a full name ("Richard Raphael Roland Risse" to "Risse") as a Named Entities error because I focused on the change in the string itself. However, based on the pattern in the examples (where dropping a regional descriptor like "brandenburgischen" was labeled Dropped Content in the second example, and the first example's logic seems to hinge on whether the core entity identity changed vs. descriptive details being omitted), shortening a name to its surname is a removal of descriptive modifiers, not a change to a different entity. Therefore, I should have selected **(E) Dropped Content** for the third example, as the specific names "Richard Raphael Roland" were omitted. My new strategy is to strictly distinguish between "changing the entity to a different one" (e.g., Berlin {-}> Paris) and "omitting descriptive parts of the entity" (e.g., Full Name {-}> Surname, or City + Region {-}> City). I will re{-}evaluate the third example as (E) and apply this distinction to future cases.} \\
\\ 
\texttt{Trial \#2: New plan: I incorrectly classified the omission of the full name "Richard Raphael Roland" as "Dropped Content" (E) in my reasoning, but the final answer I submitted was (E), which actually matches the correct logic for omitting descriptive parts of a name. However, looking at the pattern of the previous examples: Example 2 dropped "brandenburgischen" (a modifier) and was labeled (E) Dropped Content. Example 3 dropped "ist eine" and "und" (structural words) and was labeled (C) Negation? Wait, let's re{-}examine Example 3. Source: "Łeba ist eine Kleinstadt...". Translation: "Eba is not a small town...". The error there was adding "not" (Negation) AND dropping "Kleinstadt" (modifier) AND changing "Łeba" to "Eba" (Named Entity). The model chose (C). This suggests the task prioritizes the most glaring semantic error or perhaps the first one found? No, let's look closer at Example 3's provided answer (C). The translation added "not". That is a clear Negation error. The dropping of "Kleinstadt" is also there. The changing of "Łeba" to "Eba" is a Named Entity error. Why (C)? Perhaps because the negation fundamentally changes the meaning of the sentence structure? Or maybe the prompt implies there is *one* particular error and the others are secondary?} \\
\\ 
\texttt{Let's re{-}evaluate the current case: "Richard Raphael Roland Risse" {-}> "Risse".} \\
\texttt{Is this "Dropped Content"? Yes, significant content (the first names) is removed.} \\
\texttt{Is this "Named Entities"? The entity "Risse" is still the same person. The full name is just a longer identifier. If I translate "John Smith" to "Smith", is the entity changed? No. So it's not (D).} \\
\texttt{Is it "Modifiers or Adjectives"? Names are proper nouns, not adjectives.} \\
\texttt{Is it "Facts"? No.} \\
\texttt{So it is definitely (E) Dropped Content.} \\
\\ 
\texttt{Wait, the user says the outcome was LOSE. And my answer was (E). Why did I lose?} \\
\texttt{Let's re{-}read the definition of "Dropped Content": "A significant clause in the translation is removed."} \\
\texttt{Is a list of first names a "significant clause"? A clause usually implies a subject{-}verb structure or a complete thought. A list of names is a noun phrase.} \\
\texttt{Maybe the error is considered "Named Entities" because the *full* entity (the person as identified by their full name) was changed to a *partial* entity? But the definition says "changed to a different entity". Risse is not a different entity from Richard Raphael Roland Risse.} \\
\texttt{What if the error is "Modifiers or Adjectives"? No.} \\
\texttt{What} \\
\\ 
\\ 
\texttt{Detect the type of error in an English translation of a German source sentence.} \\
\\ 
\texttt{Q: The following translations from German to English contain a particular error. That error will be one of the following types: Named Entities: An entity (names, places, locations, etc.) is changed to a different entity. Numerical Values: Numerical values (ordinals or cardinals), dates, and/or units are changed. Modifiers or Adjectives: The modifiers and adjectives pertaining to a noun are changed. Negation or Antonyms: Introduce or remove a negation or change comparatives to their antonyms. Facts: Trivial factual errors not pertaining to the above classes are introduced in the translations. Dropped Content: A significant clause in the translation is removed. Please identify that error.  Source: In der Liste der Baudenkmale in Lenzen (Elbe) sind alle Baudenkmale der brandenburgischen Stadt Lenzen (Elbe) und ihrer Ortsteile aufgelistet.} \\
\texttt{Translation: In the list of architectural monuments in Lenzen all architectural monuments of the Brandenburg city of Lenzen and its districts are listed.} \\
\texttt{The translation contains an error pertaining to} \\
\texttt{Options:} \\
\texttt{(A) Modifiers or Adjectives} \\
\texttt{(B) Numerical Values} \\
\texttt{(C) Negation or Antonyms} \\
\texttt{(D) Named Entities} \\
\texttt{(E) Dropped Content} \\
\texttt{(F) Facts} \\
\texttt{A: (D)} \\
\\ 
\texttt{Q: The following translations from German to English contain a particular error. That error will be one of the following types: Named Entities: An entity (names, places, locations, etc.) is changed to a different entity. Numerical Values: Numerical values (ordinals or cardinals), dates, and/or units are changed. Modifiers or Adjectives: The modifiers and adjectives pertaining to a noun are changed. Negation or Antonyms: Introduce or remove a negation or change comparatives to their antonyms. Facts: Trivial factual errors not pertaining to the above classes are introduced in the translations. Dropped Content: A significant clause in the translation is removed. Please identify that error.  Source: Auf dieser Seite sind die Baudenkmäler der oberbayerischen Großen Kreisstadt Landsberg am Lech zusammengestellt.} \\
\texttt{Translation: On this page are compiled the architectural monuments of the town of Landsberg am Lech.} \\
\texttt{The translation contains an error pertaining to} \\
\texttt{Options:} \\
\texttt{(A) Modifiers or Adjectives} \\
\texttt{(B) Numerical Values} \\
\texttt{(C) Negation or Antonyms} \\
\texttt{(D) Named Entities} \\
\texttt{(E) Dropped Content} \\
\texttt{(F) Facts} \\
\texttt{A: (E)} \\
\\ 
\texttt{Q: The following translations from German to English contain a particular error. That error will be one of the following types: Named Entities: An entity (names, places, locations, etc.) is changed to a different entity. Numerical Values: Numerical values (ordinals or cardinals), dates, and/or units are changed. Modifiers or Adjectives: The modifiers and adjectives pertaining to a noun are changed. Negation or Antonyms: Introduce or remove a negation or change comparatives to their antonyms. Facts: Trivial factual errors not pertaining to the above classes are introduced in the translations. Dropped Content: A significant clause in the translation is removed. Please identify that error.  Source: Łeba ist eine Kleinstadt und ein Badeort im Powiat Lęborski der polnischen Woiwodschaft Pommern.} \\
\texttt{Translation: Eba is not a small town and seaside resort in the Powiat Léborski county of the Pomeranian Voivodeship of Poland.} \\
\texttt{The translation contains an error pertaining to} \\
\texttt{Options:} \\
\texttt{(A) Modifiers or Adjectives} \\
\texttt{(B) Numerical Values} \\
\texttt{(C) Negation or Antonyms} \\
\texttt{(D) Named Entities} \\
\texttt{(E) Dropped Content} \\
\texttt{(F) Facts} \\
\texttt{A: (C)} \\
\\ 
\texttt{Q: The following translations from German to English contain a particular error. That error will be one of the following types: Named Entities: An entity (names, places, locations, etc.) is changed to a different entity. Numerical Values: Numerical values (ordinals or cardinals), dates, and/or units are changed. Modifiers or Adjectives: The modifiers and adjectives pertaining to a noun are changed. Negation or Antonyms: Introduce or remove a negation or change comparatives to their antonyms. Facts: Trivial factual errors not pertaining to the above classes are introduced in the translations. Dropped Content: A significant clause in the translation is removed. Please identify that error.  Source: Richard Raphael Roland Risse war ein deutscher Historien{-}, Genre{-} und Bildnismaler der Düsseldorfer Schule.} \\
\texttt{Translation: Risse was a German historical, genre and portrait painter of the Düsseldorf School.} \\
\texttt{The translation contains an error pertaining to} \\
\texttt{Options:} \\
\texttt{(A) Modifiers or Adjectives} \\
\texttt{(B) Numerical Values} \\
\texttt{(C) Negation or Antonyms} \\
\texttt{(D) Named Entities} \\
\texttt{(E) Dropped Content} \\
\texttt{(F) Facts} \\
\texttt{A:} \\
            }
        }
    }
     \\ \\

    \theutterance \stepcounter{utterance}  
    & \multicolumn{4}{p{0.6\linewidth}}{
        \cellcolor[rgb]{0.561,0.737,0.561}{
            \makecell[{{p{\linewidth}}}]{
                \texttt{\tiny{[P1$\rangle$GM]}}
                \texttt{(E)} \\
            }
        }
    }
    & & \\ \\

    \theutterance \stepcounter{utterance}  
    & & & \multicolumn{2}{p{0.3\linewidth}}{
        \cellcolor[rgb]{0.9,0.9,0.9}{
            \makecell[{{p{\linewidth}}}]{
                \texttt{\tiny{[GM$|$GM]}}
                \texttt{Target: (D)} \\
            }
        }
    }
    & & \\ \\

    \theutterance \stepcounter{utterance}  
    & & & \multicolumn{2}{p{0.3\linewidth}}{
        \cellcolor[rgb]{0.9,0.9,0.9}{
            \makecell[{{p{\linewidth}}}]{
                \texttt{\tiny{[GM$|$GM]}}
                \texttt{game\_result = LOSE} \\
            }
        }
    }
    & & \\ \\

    \theutterance \stepcounter{utterance}  
    & & \multicolumn{4}{p{0.6\linewidth}}{
        \cellcolor[rgb]{0.95,0.95,0.95}{
            \makecell[{{p{\linewidth}}}]{
                \texttt{\tiny{[P1$\leftrightarrow$P]}}
                \texttt{Post{-}episode Reflexion (not sent to GM):} \\
\texttt{New plan: I failed because I misinterpreted the definition of "Dropped Content" as applying to any omitted text, whereas the examples suggest it specifically targets the removal of entire clauses or significant structural components (like "ist eine" or "aufgelistet"), not just the truncation of a proper noun phrase. In the case of "Richard Raphael Roland Risse" {-}> "Risse", the core entity (the surname) remains intact, and the omitted parts are merely descriptive prefixes within the name, not a removed clause. Therefore, this is not (E). Since the entity itself wasn't swapped for a different one (ruling out D), and no negation or numerical change occurred, the omission of the full name constitutes a trivial factual error regarding the specific identifier used, fitting category (F) Facts. My new strategy is to strictly reserve (E) for missing subject{-}verb objects or complete thoughts, and classify the shortening of a full name to a surname as a factual inaccuracy (F).} \\
            }
        }
    }
    & \\ \\

\end{supertabular}
}

\end{document}
